\begin{textsample}{POS Dim 3 – rs – Score 16.00 – rs/rs\_inf\_17.txt}  \label{ex:f3_pos_226}
11.5 CAMPO DE EXPERIÊNCIAS : ESPAÇOS, TEMPOS, QUANTIDADES, RELAÇÕES E TRANSFORMAÇÕES ( ET ) Neste Campo, destacam -se experiências nas quais as \textbf{crianças} falam, descrevem, narram, explicam e fazem relações, requisitos fundamentais para a construção e ampliação de saberes, como forma de fortalecimento de sua autonomia e ricas oportunidades para a construção do raciocínio lógico, de noções de espaços, tempos, quantidades, de classificação, seriação, entre outros, para chegar à percepção de relações e transformações em todas as suas vivências.
No Campo de Experiências Espaços, Tempos, Quantidades, Relações e Transformações, os \textbf{bebês} e as \textbf{crianças} \textbf{demonstram} interesse e \textbf{curiosidade} nas situações de criação de cenários, narrativas de histórias, fazem diferentes descobertas, sendo capazes de resolverem situações-problemas do cotidiano.
As \textbf{crianças}, por meio da \textbf{curiosidade} que lhes é peculiar, da indagação, da experimentação e da formulação de noções intuitivas, vão formulando questões acerca do mundo e de si mesmas.
Na busca pela ampliação do mundo físico e sociocultural das \textbf{crianças}, este Campo de Experiências, baseado no seu conceito proposto pela BNCC (2017)[10 ], envolve tanto os conhecimentos cotidianos quanto aqueles historicamente acumulados pela humanidade.
Esses conhecimentos compõem o patrimônio científico, ambiental e tecnológico, estabelecendo relações com as noções de espaço, de tempo, de quantidade, de relações e de transformações de elementos.
Diante disso, os \textbf{bebês} e as \textbf{crianças}, sejam do contexto do campo ou moradoras de zonas urbanas, estão imersas em um meio repleto de produtos da cultura.
Em sua atuação protagonista, as \textbf{crianças} buscam compreender o funcionamento das coisas que as cercam, diferenciam propriedades e características de materiais diversos, sempre fazendo indagações do “ como? ” e “ por quê? ”, buscando apropriar -se do conhecimento de maneira crítica, criativa e significativa. [ 10 ] : Espaços, tempos, quantidades, relações e transformações – As \textbf{crianças} vivem inseridas em espaços e tempos de diferentes dimensões, em um mundo constituído de fenômenos naturais e socioculturais.
Desde muito \textbf{pequenas}, elas procuram se situar em diversos espaços ( rua, bairro, cidade etc. ) e tempos ( dia e noite ; hoje, ontem e amanhã etc. ). \textbf{Demonstram} também \textbf{curiosidade} sobre o mundo físico ( seu próprio corpo, os fenômenos atmosféricos, os animais, as plantas, as transformações da natureza, os diferentes tipos de materiais e as possibilidades de sua manipulação etc. ) e o mundo sociocultural ( as relações de parentesco e sociais entre as pessoas que conhece ; como vivem e em que trabalham essas pessoas ; quais suas tradições e seus costumes ; a diversidade entre elas etc. ).
Além disso, nessas experiências e em muitas outras, as \textbf{crianças} também se deparam, \textbf{frequentemente}, com conhecimentos matemáticos ( contagem, ordenação, relações entre quantidades, dimensões, medidas, comparação de \textbf{pesos} e de comprimentos, avaliação de distâncias, reconhecimento de formas geométricas, conhecimento e reconhecimento de numerais cardinais e ordinais etc. ) que igualmente aguçam a \textbf{curiosidade}.
Portanto, a Educação \textbf{Infantil} precisa promover experiências nas quais as \textbf{crianças} possam fazer observações, \textbf{manipular} objetos, investigar e explorar seu entorno, levantar hipóteses e consultar fontes de informação para buscar respostas às suas \textbf{curiosidades} e indagações.
Assim, a instituição escolar está criando oportunidades para que as \textbf{crianças} ampliem seus conhecimentos do mundo físico e sociocultural e possam utilizá -los em seu cotidiano. ( BNCC, 2017, p. 40-41 ). 11.5.1 Direitos de Aprendizagem e Desenvolvimento no Campo de Experiências Espaços, Tempos, Quantidades, Relações e Transformações CONVIVER com \textbf{crianças} e \textbf{adultos} e com eles investigar o mundo natural e social. \textbf{BRINCAR} com materiais, objetos e elementos da natureza e de diferentes culturas e perceber a diversidade de formas, \textbf{texturas}, cheiros, cores, tamanhos, \textbf{pesos}, densidades que apresentam.
EXPLORAR características do mundo natural e social, nomeando -as, agrupando -as e ordenando -as segundo critérios relativos às noções de espaços, tempos, quantidades, relações e transformações.
PARTICIPAR de atividades de investigação de características de elementos naturais, objetos, situações, espaços, utilizando ferramentas de exploração ( bússola, lanterna, lupa ) e instrumentos de \textbf{registro} e comunicação, como máquina \textbf{fotográfica}, filmadora, gravador, projetor e computador.
EXPRESSAR suas observações, explicações e representações sobre objetos, organismos vivos, fenômenos da natureza, características do ambiente.
CONHECER -SE e construir sua identidade pessoal e cultural, reconhecendo seus interesses na relação com o mundo físico e social. 11.5.2 Objetivos de Aprendizagem e Desenvolvimento Os Objetivos de aprendizagem e desenvolvimento para os \textbf{bebês}, as \textbf{crianças} bem \textbf{pequenas} e as \textbf{crianças} \textbf{pequenas} assim são descritos na BNCC e no Referencial Curricular Gaúcho : No título seguinte, é apresentada a síntese das aprendizagens e desenvolvimento propostos pela BNCC ( 2017 ), considerando o percurso das \textbf{crianças} da Educação \textbf{Infantil} em cada um dos Campos de Experiências.
O propósito dessa síntese é tornar claro para a escola e para o professor o que é preciso organizar para garantir os direitos e as especificidades do processo educativo de todas as \textbf{crianças} nesta etapa.

% matched lemmas: adulto, bebê, brincar, criança, curiosidade, demonstrar, fotográfico, frequentemente, infantil, manipular, pequeno, peso, registro, textura
\end{textsample}
